\section{Proofs for norms}
\label{app:operation-foundations}

For a nonzero $p$-torsionfree perfectoid ring $R$, we prove two formulas
used in Subsection~\ref{subsec:graded-norms}.

\subsection{Maps induced by geometric fixed points}
\label{app:ghosts}

We compute the maps $\gamma_i$ of
Definition~\ref{def:geometric-fixed-point-maps}
for $U=\rho_p(\mathbf a)$ with $\mathbf a\in D_n^{(0)}$.

Put $L=I_0/I_0^2$, $\lambda=\mathbb C(1)$. Identify $\pi_{2a}X_R=L^a$ for $a\ge0$
by the restriction $L\xrightarrow{\sim}\pi_2X_R$ in degree $\lambda$
\cite[Example~3.2 and Remark~3.5]{SulymaNew}.
For $P=X_R,\varrho_p^*(X_R^{tC_p})$, the Euler map on the two-cell of
$B\mathbb T$ sends a weight-one class restricting to $z$ to
$c_1(\lambda)z\in H^2(B\mathbb T;\pi_2P)$.
With the same generator via $\varrho_p$, cyclotomic Frobenius induces
$\varphi:L\to I_1/I_1^2$ on these quotients and
$x\mapsto\Phi_B(1\otimes x)$ on $\pi_2$ of circle homotopy fixed points
(Subsection~\ref{subsec:period-line}, \cite[Proposition~6.2 and Remark~6.6]{BMS}).
Hence integer-degree restriction is
\[
\delta:I_0B/I_0^2B\xrightarrow{\sim}L,\qquad
[x]\longmapsto\varphi^{-1}\!\left(\operatorname{pr}_1\Phi_B(1\otimes x)\right),
\]
where $\operatorname{pr}_1:B\to B/I_1B=I_1/I_1^2$ is the first projection.
Write
\begin{equation}\label{eq:geometric-kappa}
\delta=\mathrm{id}_L\otimes\kappa,\qquad
\kappa:B/I_0B\xrightarrow{\sim}R.
\end{equation}
The change to HKR coordinates is $\delta^{-1}$ followed by the second
map in \eqref{eq:period-cotangent-reduction}.

\begin{definition}\label{def:geometric-fixed-point-coefficients}
For $0\le i<n$ and $a_i\ge0$, use $\Sigma_{n-i}$ from \eqref{eq:period-repair-data} to identify the lines
\[
\varphi^{-i*}\bigl(I_i^{a_i}E_i(U)\bigr)
\cong I_0^{a_i}B^{a_n}
\prod_{i<j<n}(J_0^{j-i})^{a_j-a_{j+1}}.
\]
Define
\[
\bar\theta_{a_i}:
\bigl(\varphi^{-i*}\bigl(I_i^{a_i}E_i(U)/I_i^{a_i+1}E_i(U)\bigr)\bigr)[1/p]
\longrightarrow (I_0/I_0^2)^{a_i}[1/p]
\]
by tensoring the identity on $(I_0/I_0^2)^{a_i}$ with the
$a_n$th tensor power of $\kappa$ from \eqref{eq:geometric-kappa} and the
$(a_j-a_{j+1})$th tensor powers of
\[
J_0^{j-i}/I_0J_0^{j-i}\longrightarrow R,
\qquad [a]\longmapsto\theta(a)
\quad(i<j<n).
\]
For negative powers, use the duals of the inverse maps after
inverting $p$.
\end{definition}

\begin{proposition}\label{prop:geometric-fixed-point-coordinates}
Let $x=(x_0,\ldots,x_n)\in\ker\Delta_U$.
For $0\le i<n$, if $a_i<0$, then $\gamma_i(x)=0$.
If $a_i\ge0$, let $\overline{x_i}$ be the image of $x_i$ in
$I_i^{a_i}E_i(U)/I_i^{a_i+1}E_i(U)$.
Then
\begin{equation}\label{eq:ghost-line-conversion}
\gamma_i(x)
=p^{-\sum_{i<j<n}(a_j-a_n)}
\bar\theta_{a_i}\bigl(\varphi^{-i*}\overline{x_i}\bigr)
\in L^{a_i}.
\end{equation}
\[
p^{-\sum_{i<j<n}(a_j-a_n)}\bar\theta_{a_i}:
\varphi^{-i*}\bigl(I_i^{a_i}E_i(U)/I_i^{a_i+1}E_i(U)\bigr)
\xrightarrow[\ R\text{-linear}\ ]{\sim}(I_0/I_0^2)^{a_i}.
\]
At the terminal vertex,
\[
\gamma_n(x)=\delta^{\otimes a_n}\bigl(\varphi^{-n}(x_n)\bigr)
\in L^{a_n}\qquad(a_n\ge0),
\]
where $A_n(U)=\varphi^{n*}\bigl((I_0B/I_0^2B)^{\otimes_R a_n}\bigr)$.
If $a_n<0$, then $\gamma_n(x)=0$.
\end{proposition}

\begin{proof}
By \cite[Theorem~6.1]{BMS}, it suffices to consider $a_i\ge0$.
By \cite[Lemma~2.2.8]{BhattLurie} and $p$-torsionfreeness of $R$,
\[
J_0^k/I_0J_0^k\xrightarrow[\sim]{\ \theta/p^k\ }R
\quad(k\ge1).
\]
Taking tensor powers, with
$\sum_{i<j<n}(j-i)(a_j-a_{j+1})=\sum_{i<j<n}(a_j-a_n)$,
and then using $\kappa$ proves the isomorphism.

For $k\ge1$, pullback along $z\mapsto z^{p^k}$ induces
$I_0J_0^k\hookrightarrow I_0$ in the Euler coordinates
of \cite[Notation~2.11 and Example~3.2]{SulymaNew}. Its underlying
degree is $p^k$, so $p^k\res_e^{\mathbb T}(z)=[z]\in L$ for $z\in I_0J_0^k$.
Thus restriction in degrees $\mathbb C(p^k)-\lambda$ and
$\mathbb C(0)-\lambda$ is $\theta/p^k$ and $\kappa$, respectively.
Tensor products and duals give the restriction for $U^{C_{p^i}}$
\cite[Remark~3.1 and Examples~3.3--3.4]{SulymaNew}.

For the circle lift of Definition~\ref{def:app-local-spectra},
Proposition~\ref{prop:app-finite-comparison} identifies restriction
to $C_{p^{n-i}}$ with the quotient onto $\varphi^{-i*}A_i(U)$.
Since $\res_e^{\mathbb T}
=\res_e^{C_{p^{n-i}}}\res_{C_{p^{n-i}}}^{\mathbb T}$,
composition with the $i$th vertex of \eqref{eq:tate-compression}
gives \eqref{eq:ghost-line-conversion}.
\end{proof}

\subsection{Proof of \eqref{eq:full-boundary-norm}}
\label{app:full-cut-norms}

For $H=C_{p^h}<G=C_{p^n}$ and $0\le i<h$, we compute the norm of a
boundary class. We show that $\partial_i$ respects the indexed
multiplication for $G/H$.

\begin{proposition}\label{prop:boundary-indexed-multiplications}
Let $H=C_{p^h}<G=C_{p^n}$ and $0\le i<h$.
For the indexing system~\cite[Definition~3.22]{BHOperadic} generated by $G/H$, the Tate compression
diagram truncated at $X_R^{\Phi_{i+1}}$ is a diagram of algebras, and
\[
\partial_i:\Omega X_R^{t_i}\longrightarrow X_R
\]
is a map of nonunital algebras.
Consequently, \eqref{eq:full-boundary-norm} holds.
\end{proposition}

\begin{proof}
For $-1\le r\le i$, we have $C_{p^{r+1}}\le H$.
By \cite[Theorem~6.1 and Corollary~6.5]{BHNorm}, the functors
$\widetilde E\mathcal F_r\wedge(-)$ preserve these algebras,
and the localization maps are maps of algebras.
The $K$-equivariant diagonals $E\mathcal F_r\to\prod_T E\mathcal F_r$,
for all admissible finite $K$-sets $T$ and $K\le G$,
make $F((E\mathcal F_r)_+,-)$ a functor on these algebras,
and the completion maps are maps of algebras.

By Proposition~\ref{prop:compression}, the $\mathcal F_r$-completion
maps of $X_R^{t_{r-1}}$ are equivalences of these algebras
for $1\le r\le i$.
Naturality of completion gives the factorizations of
Proposition~\ref{prop:compression} in the same category.
Iterated isotropy separation gives
\begin{equation}\label{eq:full-cut-square}
X_R\simeq
\holim\bigl(X_R^{h_0}\to X_R^{t_0}\leftarrow\cdots\to
X_R^{t_{i-1}}\leftarrow X_R^{h_i}\bigr)
\mathop{\times}_{X_R^{t_i}}X_R^{\Phi_{i+1}},
\end{equation}
where the homotopy limit is $X_R^{h_0}$ for $i=0$.

For a homotopy pullback $B\times_D C$ of these algebras,
the zero maps to $B$ and $C$ give a map of nonunital algebras
\[
\Omega D=0\mathop{\times}_D0\longrightarrow B\times_D C.
\]
Limits of nonunital algebras are computed on spectra,
so the map for \eqref{eq:full-cut-square} is $\partial_i$.
The indexed multiplication on $\Omega D\simeq S^{-1}\wedge D$ is
\[
S^{-\mathbb R[G/H]}\wedge N_H^G\res_H^G D
\longrightarrow S^{-\mathbb R[G/H]}\wedge D
\xrightarrow{\ a_H^G\ }S^{-1}\wedge D.
\]
$a_H^G$ is dual to the one-point compactification of
\[
\mathbb R\hookrightarrow\mathbb R[G/H],
\qquad t\longmapsto(t)_{gH\in G/H}.
\]
Then taking homotopy groups gives \eqref{eq:full-boundary-norm}
(cf.\ \cite[Theorem~3.1.1]{Balderrama} for the Euler factor).
\end{proof}

